\section{Martingale：公平游戏与停时}
\label{sec:06-Random-Processes/04-Applications/02}

赌博直觉里有一个朴素想法：如果游戏公平，你再怎么决定何时停手，期望收益总归是零。这东西到底对不对？如果对，在什么条件下对？为了精确回答，需要一整套语言。

设 \(\mathcal F_n\) 表示到时刻 \(n\) 为止你掌握的全部信息——不只是当前数值，还包括过去的整条轨迹。正式说法是，\(\mathcal F_n\) 是一族随 \(n\) 递增的 \(\sigma\)-代数（filtration）。一个随机过程 \(M_n\) 称为关于这组信息\textbf{适应}的，是指 \(M_n\) 的值在时刻 \(n\) 已被 \(\mathcal F_n\) 完全确定——你不能让今天的状态依赖于明天的信息。

如果还满足

\[
\E[|M_n|]<\infty,
\]

以及

\[
\E[M_{n+1}\mid\mathcal F_n]=M_n,
\]

就称 \(M_n\) 是一个\textbf{Martingale}。

翻译回直觉：你已经知道了到时刻 \(n\) 为止的全部历史，在此条件下，下一步的期望值就是当前值。没有可预见的平均漂移——无论是涨是跌，都均值为零。

但\textbf{没有漂移不等于不会动}。Martingale 路径完全可以剧烈起伏，方差可以越来越大。公平硬币累加的位置就是 Martingale，可是走了一万步之后，位置分布的宽度大约是 100 步，绝非安稳地停在原点。

\subsection{公平游走：最简单的 Martingale}

令

\[
S_n=Z_1+\cdots+Z_n,
\]

其中增量 \(Z_{n+1}\) 与过去独立，均值零。信息族取为

\[
\mathcal F_n=\sigma(Z_1,\ldots,Z_n),
\]

即``由前 \(n\) 步增量可回答的所有事件''构成的集合。那么

\[
\E[S_{n+1}\mid\mathcal F_n]
=S_n+\E[Z_{n+1}\mid\mathcal F_n]
=S_n.
\]

增量独立于过去，且均值为零，所以条件期望把增量吞掉，剩下 \(S_n\)。公平游走满足 Martingale 性，和直觉吻合。

如果增量均值为正，则

\[
\E[M_{n+1}\mid\mathcal F_n]\ge M_n,
\]

称为 submartingale——趋势向上。均值为负则为 supermartingale——趋势向下。Martingale 是两者在零漂移处的重合点。

\subsection{条件期望生来就是 Martingale}

有一个比赌博更深刻的构造。取任意可积随机变量 \(X\)，定义

\[
M_n=\E[X\mid\mathcal F_n].
\]

\(M_n\) 是随着信息不断增多时，你对最终未知量 \(X\) 的当前最佳预测。由塔式性质，

\[
\E[M_{n+1}\mid\mathcal F_n]
=\E[\E[X\mid\mathcal F_{n+1}]
\mid\mathcal F_n]
=\E[X\mid\mathcal F_n]
=M_n.
\]

所以``逐步揭示信息后的条件预测''天然构成 Martingale，不管你预测的是什么东西。这个结构把 Martingale 从赌场的具体场景中抽离出来：它的本质是信息过滤与条件期望的一致性，而不是赌局本身的公平性。

设想 \(X\) 是你今天通勤最终花费的总分钟数，\(\mathcal F_n\) 包含你走到第 \(n\) 个路口时已经知道的用时、当前位置和沿途交通消息。那么

\[
M_n=\E[X\mid\mathcal F_n]
\]

就是你对总用时的不断更新的预测。走到第三个路口时前方突然拥堵，\(M_3\) 上调；过了拥堵段道路畅通，\(M_7\) 下调。Martingale 性允许预测值随新信息而变动，但它不允诺任何方向：如果经过某个路口后预测总是系统性上调，说明当时掌握的规律还未穷尽可得信息，或者模型漏掉了某些可预测的结构。

同样的框架换一个名字就能换一个场景：把最终通勤时长换成日终总订单量，把路口信息换成逐小时订单、流量和活动数据，日终预测过程本身仍是 Martingale。

\subsection{补偿把漂移摘掉}

如果随机游走单步均值为 \(\mu\neq 0\)，\(S_n\) 本身有漂移，不是 Martingale。但你可以手动摘掉漂移：

\[
M_n=S_n-n\mu
\]

满足 Martingale 性——你把可预测的趋势分量从观测量中减去，留下的就是不可预测的创新。对 Poisson 过程，

\[
N(t)-\lambda t
\]

也是连续时间 Martingale：计数减掉期望速率，结果的条件均值为零。

``观测量减去可预测趋势，得到 Martingale''是时间序列分解、计数过程分析和统计估计中的共同线索。ARIMA 模型里的差分操作也好、Doob--Meyer 分解也好，都是在做同一件事：从过程中分离出可预测部分，剩下的就是 Martingale。

\subsection{停时：你不能偷看未来}

设想你有一个策略，看到过程涨到某个水平就停止，或者跌破某个水平就停止。停时 \(\tau\) 的规则是：在时刻 \(n\) 决定是否停止，只能使用 \(\mathcal F_n\) 中已有的信息。不能偷看 \(n+1\) 及其后数据再做决定。形式条件是 \(\{\tau=n\}\in\mathcal F_n\) 对所有 \(n\) 成立。

可选停止定理说：在适当条件下，

\[
\E[M_\tau]=\E[M_0].
\]

不管你的停止策略多么精巧，只要用的是当前及过去的信息，停止时刻的期望值不会偏离初始值。合理——游戏本身公平，你无权以策略制造不公平。

但``适当条件''四个字不能省略。过去有一群人试图用加倍下注策略（输一局就加倍下注，直到赢回为止）来保证最终净赢一元。表面上，只要最终赢一局，就能收复所有损失再加一元。矛盾在哪？

加倍下注需要无限信用来支撑几何增长的赌注，也需要无限时间来保证最终会赢——但赌局可能在你输光或被赶走之前还没有赢到那一局。这个策略里，下注金额无界，停止时间期望可能无限，可选停止定理的条件不成立。现实中的信用限额和赌注上限把无限理想砍断，Martingale 结论就此恢复。

常见的充分条件包括：

\begin{itemize}
\tightlist
\item
  \(\tau\) 有界——赌局有固定上限场次；
\item
  \(M_n\) 增量有界且 \(\E[\tau]<\infty\)——每步波动受限且平均等待有限；
\item
  停止后的变量族一致可积——分布尾部不逃逸。
\end{itemize}

\subsection{停时算破产概率：换个工具，同一个结果}

公平游走在 \((0,N)\) 吸收，从 \(i\) 出发。令

\[
\tau=\inf\{n:S_n\in\{0,N\}\}.
\]

在有限区间中，\(\tau\) 几乎必然有限——上一单元首步分析还给出了 \(\E_i[\tau]=i(N-i)<\infty\)。停止前 \(S_n\) 始终落在 \([0,N]\) 内，因此停止过程 \(S_{n\wedge\tau}\)（\(a\wedge b\) 取两者较小值）有界。可选停止的充分条件成立：

\[
\E_i[S_\tau]=\E_i[S_0]=i.
\]

设 \(h=P_i(S_\tau=N)\)（即从 \(i\) 出发先到 \(N\) 的概率），则

\[
\E_i[S_\tau]=Nh+0\cdot(1-h),
\]

所以

\[
h=\frac iN.
\]

同样一个赌徒破产概率，既可以用首步分析的局部递推一层层剥出来，也可以用 Martingale 停时的全局守恒量一步到位。前者在细节中递推，后者在整体中守恒。两个视角互补，不是重复劳动。

\subsection{Martingale 不保证的事}

Martingale 性只控制条件期望为零，不控制其他任何性质：

\begin{itemize}
\tightlist
\item
  增量可以有条件异方差——波动时大时小，只要正负方向均值抵消即满足 Martingale；
\item
  增量可以有复杂依赖——ARCH、GARCH、随机波动率模型都可以是 Martingale 差序列；
\item
  \(M_n\) 的分布可以随时间扩散，绝非平稳过程——方差可以随 \(n\) 增长；
\item
  路径可以有跳跃、平台期、振荡，只是没有可预测的平均趋势。
\end{itemize}

Martingale 控制的是``没有可预测的平均漂移''，不是``没有风险''，不是``未来独立于过去''，也不是``每步大小相同''。把它当作``安全''或``平稳''的同义词，是对它的最严重误解。
